\section{Related Works}

A review of literature was conducted to explore the available research in the area of 
classification of bird species, fine-grained visual recognition, and citizen science 
platforms. The selected literature was critically reviewed and analyzed based on their 
choice of architecture of machine learning algorithm, approaches for dataset creation 
and evaluation, and performance results.

Traditional bird classification systems were typically based on classic machine learning
techniques, such as SVM and Random Forest classifiers. Systems like these would depend
on manually designed features like color histograms, shape, and texture descriptors.
However, these techniques would only achieve limited success when dealing with real-world
image databases which have varying illumination, background, image quality,
and bird poses \cite{kambali2024automated, lolge2025evolution}. Recent studies
have shown that CNN architectures can significantly improve the results compared to
traditional machine learning techniques. This is due to their ability to automatically extract hierarchical
visual features from the image \cite{kambali2024automated}.

Bird species classification can be considered as a Fine-Grained Visual Classification
(FGVC) problem because bird species can share very similar visual features.
When working with FGVC tasks, models must distinguish subtle differences such as
feather patterns, beak structure, wing shape, and color variations. Chen et al. emphasize
that successful FGVC systems require models capable of learning discriminative local
features and visual details \cite{chen2024fetfgvc}. Modern deep learning architectures,
especially CNN- and Transformer-based models, can provide the performance necessary for
handling classification tasks like this.

There are several studies conducted that analyze the performance of different CNN models for bird species classification.
For example, the study by Bilgin et al. compared three different architectures such as VGG16, ResNet50, and DenseNet. 
DenseNet and ResNet were observed to perform better than VGG16 due to better feature propagation and gradient flow \cite{bilgin2022bird}.
Another example is the study conducted by Tang, which compares VGG, ResNet, InceptionV3, and MobileNet and shows that ResNet and Inception 
achieve better F1 scores, and MobileNet gives better computational efficiency \cite{tang2025comparative}.

Apart from the accuracy of classification, recent developments have also focused on
inference speed and model size \cite{lolge2025evolution, tang2025comparative}.
Models such as MobileNet proved useful where the deployment of the model is done
on websites and mobile applications. However, dense and residual neural network models
tend to perform more effectively but at the cost of increased memory and computational
requirements \cite{bilgin2022bird}. In our case, we were constrained by this factor,
since the model had to be deployed as part of a live website.

The literature survey also included public datasets with images of birds and citizen science platforms. 
Large scale biodiversity datasets are available, for example the Global Biodiversity Information Facility (GBIF) \cite{gbif2025}. 
These resources offer free datasets that contain records collected from all over the world.

Citizen science has transformed biodiversity monitoring through decentralized data collection from large communities of volunteers.
The authors Fraisl et al. argue that the ability of citizen science to gather data on a large scale makes it an important part of ecological research \cite{fraisl2022citizen}.
Bonney et al. argue similarly that citizen science projects can increase public engagement and awareness while collecting valuable environmental data \cite{bonney2015citizen}. 
However, citizen science data can pose challenges such as observer bias, class imbalance, inconsistent image quality, and incorrect labels \cite{johnston2022outstanding}.
Automated verification and validation techniques are increasingly being used to enhance the quality and reliability of citizen science data \cite{baker2021verification}.

The results from the literature review shaped many of the design choices made in the project.
The results were used in the selection of architectures, the criteria for the performance assessment, and the data pre-processing steps.
Special attention was paid to issues such as class imbalance, specie similarity, and varying image quality

\clearpage